\paragraph{C.3y (Proved/refuted): an explicit multiplicative ordinary divisor quotient in every characteristic.}
For every algebraically closed field $k$, there is an explicitly
specified excellent normal strictly henselian two-dimensional local
$k$-algebra $R$, with regular puncture $X$, such that
\begin{equation}\label{eq:C3y-answer}
 \operatorname{Pic}(X)_{\mathrm{alg}}=0,
 \qquad
 B^1(X)=\operatorname{Pic}(X)
       \simeq k^\times\oplus(\mathbb Z/4\mathbb Z)^2.
\end{equation}
The last isomorphism is one of abstract abelian groups and is
noncanonical; the injection of $k^\times$ below is specified by
gluing on the exceptional triangle. In particular $B^1(X)$ is
infinite in every characteristic. Its rationalization is zero
when $k=\overline{\mathbb F}_p$, and is infinite dimensional over
$\mathbb Q$ in characteristic zero and whenever $k$ has positive
transcendence degree in positive characteristic. Algebraic
equivalence in \eqref{eq:C3y-answer} is ordinary equivalence by
actual product families, not membership in a local Picard
functor's identity component.

Here are equations defining the ring, including the exceptional
characteristic choice:
\[
 a=\begin{cases}2,&\operatorname{char}k=229,\\1,&\text{otherwise},\end{cases}
 \qquad
 D=xyz,\qquad
 G=x^4+y^4+z^4+a(x^3y+y^3z+z^3x),
\]
\begin{equation}\label{eq:C3y-ring}
 R=\left(k[x,y,z]/(xyz-G)\right)_{(x,y,z)}^{\mathrm{sh}},
 \qquad
 X=\operatorname{Spec}R\setminus\{\mathfrak m\}.
\end{equation}
No classification of cusp singularities and no contraction
theorem for algebraic spaces is needed for this construction.

\begin{proof}
\emph{1. Twelve transverse points and an explicit projective resolution.}
Let $D_0\subset\mathbb P^2_k$ be the coordinate triangle $xyz=0$.
The quartic $G$ is nonzero at all three vertices. On each side,
after choosing the cyclic affine coordinate with the two
vertices at zero and infinity, its restriction is
\[
                       g_a(t)=t^4+a t^3+1.
\]
The resultant with its derivative is
\[
                \operatorname{Res}(g_a,g_a')=256-27a^4.
\]
For $a=1$ this is $229$, which is nonzero in every allowed
characteristic except the explicitly excluded characteristic
$229$. There $a=2$ gives $-176\equiv53\pmod{229}$.
Thus $D_0$ and $G=0$ meet transversely in twelve distinct
points, four on each side and none at a vertex. This assertion
also holds in characteristics $2$ and $3$. Exact integer
Bezout identities, valid before reduction, are
\begin{align*}
 229&=(48t^2+229)g_1
       +(-12t^3-3t^2-55t-16)g_1',\\
 -176&=(192t^2-176)g_2
       +(-48t^3-24t^2+80t-32)g_2'.
\end{align*}

Blow up these twelve reduced points in $\mathbb P^2_k$, obtaining
a smooth projective surface $S$. Let $E_1,E_2,E_3$ be the strict
transforms of the three sides, and put $E=E_1+E_2+E_3$.
Each $E_i\simeq\mathbb P^1_k$ has self-intersection $-3$,
and each pair of different components meets transversely
in its unchanged triangle vertex. Consequently
\begin{equation}\label{eq:C3y-matrix}
 M=(E_i\cdot E_j)_{i,j}
 =\begin{pmatrix}-3&1&1\\1&-3&1\\1&1&-3\end{pmatrix},
 \qquad \det M=-16,
 \qquad \deg\mathcal O_E(-E)|_{E_i}=1.
\end{equation}

The four quartics $G,xD,yD,zD$ have base ideal $(G,D)$
as an ideal sheaf on $\mathbb P^2$. Its base points are
exactly the twelve simple points just blown up, so they
define a morphism
\[
 S\longrightarrow W\subset\mathbb P^3_k,
 \qquad [G:xD:yD:zD],
\]
where, in coordinates $(u_0:u_1:u_2:u_3)$,
\begin{equation}\label{eq:C3y-projective-image}
             W=V_+\bigl(u_0u_1u_2u_3-G(u_1,u_2,u_3)\bigr).
\end{equation}
The equation follows from homogeneity:
$G(xD,yD,zD)=D^4G(x,y,z)$ and
$(xD)(yD)(zD)=D^4$. It is irreducible, since it is
primitive of degree one in $u_0$ and
$\gcd(u_1u_2u_3,G(u_1,u_2,u_3))=1$.

This morphism contracts exactly $E$ to
$v=(1:0:0:0)$ and is an isomorphism
\begin{equation}\label{eq:C3y-resolution-open}
                           S\setminus E\simeq W\setminus\{v\}.
\end{equation}
Here is a chart verification, including the exceptional curves
of the twelve blowups. On the chart $u_i\ne0$ of the right-hand
side, the ratios $(u_1:u_2:u_3)$ give a point of $\mathbb P^2$
in its corresponding coordinate chart. After putting its
$i$th coordinate equal to one, the remaining equation is
$G=tD$, where $t=u_0/u_i$. This is precisely the $D$-chart
of the blowup of the ideal $(D,G)$. These charts exhaust
$S\setminus E$: at a base point the omitted $G$-chart
has coordinate $D/G$, and its zero locus is the strict
transform of the triangle. Away from the base points the
same assertion is immediate. Thus there are no additional
contracted curves or unidentified puncture branches.
The image is all of the integral surface $W$, since the
morphism is proper and has the dense open image just
identified. The hypersurface $W$ is Cohen--Macaulay and
regular away from $v$ by \eqref{eq:C3y-resolution-open};
it is therefore normal. Its chart $u_0=1$ is exactly
the affine equation in \eqref{eq:C3y-ring}.

Base change $S\to W$ to $\operatorname{Spec}R$, and call
the resulting regular integral proper surface $Z$.
Strict henselization is ind-etale over the local algebraic
model; it preserves regularity of this resolution and the
specified reduced exceptional triangle. The open complement
of $E$ is $X$. Moreover
\begin{equation}\label{eq:C3y-maximal-ideal}
                        \mathfrak m\mathcal O_Z=\mathcal O_Z(-E).
\end{equation}
Indeed, near $E$, the section $G$ representing $u_0$ is
nonvanishing in the resolved linear system. The three
functions $u_i/u_0$ are $(D/G)x,(D/G)y,(D/G)z$ in local
trivializations, and $x,y,z$ generate the trivialized
tautological line bundle. The ratio $D/G$ has a simple
zero on each strict triangle component; the common
exceptional factors at the twelve base points cancel.
Equivalently, its blowup $G$-chart is exactly the local
coordinate $D/G$ cutting out the strict triangle. At
a triangle node the equation is the product of the two
branch parameters, which cuts out the reduced Cartier
divisor $E$. This proves \eqref{eq:C3y-maximal-ideal}
scheme-theoretically.

\emph{2. The formal Picard group is exactly that of the triangle.}
Let $E_n\subset Z$ be defined by $\mathcal O_Z(-nE)$,
so $E_1=E$. For every $n\ge1$,
\begin{equation}\label{eq:C3y-coherent-vanishing}
             H^1(E,\mathcal O_E(-nE))=0.
\end{equation}
To see this directly, normalize the triangle. The line bundle
restricts to $\mathcal O_{\mathbb P^1}(n)$ on each of its
three components. Each component has two attaching points,
and, for $n\ge1$, evaluation of its sections at those two
points is surjective. In the normalization exact sequence
the map from the three spaces of sections to the three
node fibres is therefore surjective. The component
$H^1$ groups vanish, proving \eqref{eq:C3y-coherent-vanishing}.
Also $H^2$ of every coherent sheaf on this curve vanishes.

Apply these facts to the square-zero unit sequence
\[
 1\longrightarrow\mathcal O_E(-nE)
 \longrightarrow\mathcal O_{E_{n+1}}^\times
 \longrightarrow\mathcal O_{E_n}^\times\longrightarrow1.
\]
Restriction gives isomorphisms
$\operatorname{Pic}(E_{n+1})\simeq\operatorname{Pic}(E_n)$,
and the maps on global units are surjective. Hence compatible
formal line bundles, including their isomorphisms, give
exactly $\operatorname{Pic}(E)$; there is no residual inverse-limit
term from units. Write $\widehat R$ for the maximal-ideal
completion and $\widehat Z=Z\times_R\widehat R$.
Proper formal existence and full faithfulness, applied using
\eqref{eq:C3y-maximal-ideal}, now give
\begin{equation}\label{eq:C3y-formal-pic}
                         \operatorname{Pic}(\widehat Z)\simeq\operatorname{Pic}(E).
\end{equation}
The exact primary theorem is Grothendieck existence for a
proper scheme over a complete noetherian ring
\citep[Theorem 30.27.1, Tag 088E]{rootStacksExistence}.
A formal invertible sheaf algebraizes to an invertible
sheaf: local freeness holds near the closed fibre by
Nakayama, and the non-locally-free closed locus has closed
image under the proper morphism, so cannot be nonempty
while avoiding the closed point of the local base.

\emph{3. Passage to the actual puncture and computation of the group.}
Excellence makes $R\to\widehat R$ regular, so $\widehat Z$
is regular. Every line bundle on
$\widehat X=\operatorname{Spec}\widehat R\setminus\{\widehat{\mathfrak m}\}$
extends to $\widehat Z$ by Cartier closure. The kernel
of restriction consists exactly of the divisors supported
on $E_1,E_2,E_3$. Thus \eqref{eq:C3y-formal-pic} gives
\[
 \operatorname{Pic}(\widehat X)
  \simeq\operatorname{Pic}(E)\big/
        \langle\mathcal O_Z(E_1)|_E,
                  \mathcal O_Z(E_2)|_E,
                  \mathcal O_Z(E_3)|_E\rangle.
\]
The Picard comparison for a noetherian henselian pair and
its completion identifies this with $\operatorname{Pic}(X)$
\citep[Section III, unnumbered Picard corollary, p.~573]{D_Elkik1973}.
Here the pair is $(R,\mathfrak m)$, and its complementary
open is exactly the puncture. Thus this calculation gives
actual line bundles on the strict henselization, not only
formal line bundles.

Normalization and gluing at the three nodes give an exact
sequence of abelian groups
\begin{equation}\label{eq:C3y-triangle-pic}
 0\longrightarrow k^\times\longrightarrow
 \operatorname{Pic}(E)\xrightarrow{\deg}
 \mathbb Z^3\longrightarrow0.
\end{equation}
Explicitly, trivial line bundles on the three components
are glued by scalars at the three nodes. Changing the
three component trivializations removes two independent
scalars, and the oriented product of the three gluing
scalars is the remaining element of $k^\times$.
For arbitrary specified component degrees, gluing the
three component line bundles similarly shows surjectivity
onto $\mathbb Z^3$. The degrees of the three exceptional
divisor classes are the columns of $M$ in
\eqref{eq:C3y-matrix}. Since $\det M\ne0$, their lattice
has zero intersection with the displayed $k^\times$.
The Smith form of $M$ is $(1,4,4)$. Consequently
\begin{equation}\label{eq:C3y-puncture-extension}
 0\longrightarrow k^\times\longrightarrow
 \operatorname{Pic}(X)\longrightarrow
 (\mathbb Z/4\mathbb Z)^2\longrightarrow0.
\end{equation}
The group $k^\times$ is divisible because $k$ is
algebraically closed. Divisible abelian groups are
injective $\mathbb Z$-modules, so this sequence splits
as abstract abelian groups. Its multiplicative subgroup
has a concrete specification: for $\lambda\in k^\times$,
glue the three trivial component bundles with scalars
$(\lambda,1,1)$ in a fixed cyclic orientation, extend
uniquely through \eqref{eq:C3y-formal-pic}, restrict to
$\widehat X$, and use Elkik's bijection to obtain the
actual class in $\operatorname{Pic}(X)$.

\emph{4. Ordinary algebraic equivalence is zero.}
Let $T$ be a smooth connected parameter curve and let
$\mathcal M$ be an actual line bundle on $X\times T$.
Complete $T$ to a smooth projective curve $C$.
The scheme $Z\times C$ is regular and noetherian, so
Cartier closure extends $\mathcal M$ to a line bundle
on $Z\times C$. Its restriction to $E\times C$ has
constant isomorphism class on the fibres over $C$.
One can verify this without any local Picard representability:
on each $E_i\times C\simeq\mathbb P^1\times C$ the
line bundle has the form
\[
 \operatorname{pr}_{E_i}^*\mathcal O_{\mathbb P^1}(d_i)
       \otimes\operatorname{pr}_C^*N_i
\]
with constant integer $d_i$ and $N_i\in\operatorname{Pic}(C)$.
Gluing at a node identifies the corresponding $N_i$ and
$N_j$, after fixed one-dimensional fibre factors are
chosen. Make these identifications along two edges of
the triangle. The last gluing is an automorphism of
one line bundle on $C$, hence multiplication by an
element of $\Gamma(C,\mathcal O_C^\times)=k^\times$.
All component degrees and the remaining gluing scalar
are therefore constant on $C$. This is the explicit
proper-curve-to-$\mathbb G_m$ constancy argument.

For $t_0,t_1\in T(k)$, the quotient of the two extended
fibre line bundles is trivial on $E$. It is trivial on
$\widehat Z$ by \eqref{eq:C3y-formal-pic}, hence on
$\widehat X$, and then on $X$ by Elkik's injectivity.
Every ordinary algebraic-equivalence generator is thus
zero. This proves $\operatorname{Pic}(X)_{\mathrm{alg}}=0$
integrally and completes \eqref{eq:C3y-answer}.

\emph{5. Precise finiteness consequences.}
Every algebraically closed field has infinite multiplicative
group, so $B^1(X)$ is infinite in every case. If
$k=\overline{\mathbb F}_p$, this multiplicative group is
torsion, and \eqref{eq:C3y-puncture-extension} shows
$B^1(X)_{\mathbb Q}=0$. In positive characteristic with
$k\ne\overline{\mathbb F}_p$, choose $t\in k$
transcendental over $k_0=\overline{\mathbb F}_p$.
The classes of $t-c$, $c\in k_0$, are linearly independent
in $k^\times\otimes\mathbb Q$: after clearing a putative
rational relation, a finite product
$\prod_c(t-c)^{n_c}$ would be a root of unity, hence a
constant in $k_0$. Its divisor on $\mathbb P^1_{k_0}$
forces every $n_c=0$. In characteristic zero, the
distinct positive rational primes give the analogous
infinite independent set by unique factorization in
$\mathbb Q^\times$. The injection in
\eqref{eq:C3y-puncture-extension} preserves this
independence after tensoring with $\mathbb Q$.
\end{proof}

The exact script
\texttt{code/verify\_multiplicative\_divisor\_quotient.py}
checks the two integer Bezout certificates, the universal
resultant and projective-image identities, the intersection
matrix and its Smith invariants, and the component evaluation
ranks for the finite sample $1\le n\le12$. Its elapsed
runtime and exact output are saved. The all-$n$ vanishing,
formal existence, strict-local descent, and infinite
multiplicative independence are proved above rather than
inferred from those computations.
